\documentclass{article}

\usepackage[T1]{fontenc}

\usepackage{booktabs,graphicx,siunitx}

\begin{document}

\section{Distribution \& normality}

\begin{tabular}{lllllll}
\toprule
Variable & Test & Statistic & N & p & Skew & Kurtosis (excess) \\
\midrule
exam\_score & Shapiro-Wilk & W 0.99 & 240 & .158 & $-$0.18 & $-$0.53 \\
exam\_score & K--S (Lilliefors) & D 0.04 & 240 & .608 & $-$0.18 & $-$0.53 \\
absences & Shapiro-Wilk & W 0.86 & 240 & $<$.001 & 1.50 & 2.43 \\
absences & K--S (Lilliefors) & D 0.17 & 240 & $<$.001 & 1.50 & 2.43 \\
\bottomrule
\end{tabular}

\includegraphics{figures/01_distribution-normality/figure_histogram_exam_score.png}

\includegraphics{figures/01_distribution-normality/figure_qq_exam_score.png}

\includegraphics{figures/01_distribution-normality/figure_histogram_absences.png}

\includegraphics{figures/01_distribution-normality/figure_qq_absences.png}

A small p (below .05) suggests the data depart from normality. With large samples these tests over-flag trivial deviations; with small samples they have low power, so a p $>$ .05 does not prove normality - it only means no departure was detected. Either way, judge normality mainly from the Q--Q plot - points hugging the diagonal indicate normality.

exam\_score: Normality was assessed with the Shapiro-Wilk test, W = .99, p = .158. absences: Normality was assessed with the Shapiro-Wilk test, W = .86, p $<$ .001.

Statistical basis: Shapiro-Wilk test for normality (3-5000 cases) (Shapiro, S. S., Wilk, M. B. (1965). "An analysis of variance test for normality (complete samples)." Biometrika, 52(3/4), 591-611.); Additional normality tests (Lilliefors and related) (Gross J, Ligges U (2015). nortest: Tests for Normality. R package.).

\end{document}